%% kapa: exam=1 round=22 year=2011 period=1 form=A subject=economics subject_ko=경제원론 seq=26 number=66
%% source: Q-Net 감정평가사 기출문제 (cst00309) artlSeq=5036790 file=extracted/1008/5036790/2011년 제22회 감정평가사 1차시험 기출문제/1교시.hwp
%% conversion: kordoc 4.24.51 -> tools/gen1.py (auto); review=proofread-sonnet; flags=none
\begin{question}{66}
독점시장에 관한 설명으로 옳은 것은? (단, 독점기업은 이윤을 극대화, 수요곡선은 우하향하는 직선)
\begin{choices}
  \item 독점기업은 시장수요곡선의 가격탄력성이 1보다 큰 구간에서 재화를 생산한다.
  \item 가격과 한계비용이 일치하는 점에서 균형이 발생한다.
  \item 단기적으로 균형에서 가격이 평균비용보다 낮으면 이익이 발생한다.
  \item 공급곡선이 존재한다.
  \item 독점기업이 직면하는 한계수입곡선은 우상향한다.
\end{choices}
\end{question}
\answer{1}
